
\lesson{26/10}
\textbf{Note} Since from Quantum Mechanics spin-$\frac{1}{2}$ particles obey
the Fermi-Dirac statistic, we expect anti-symmetry for the Dirac field and this
is provided by \textit{anti}commutators.

As before, the fields are expanded in their F.Ts:
$$\psi(x^\mu) = \int \frac{\dd^3p}{(2\pi)^{\frac{3}{2}}\sqrt{2E_p}}
\sum_{s=1}^{2}\left[u_s(\mathbf{p})a_s(\mathbf{p})e^{-ipx} +
   v_s(\mathbf{p}) b_s^\dagger(\mathbf{p})e^{ipx}
\right]
$$
where $px = P_\mu x^\mu$, $P^0=E_p$ and $u_s,v_s$ are the spinors defined earlier
in the solution of the Dirac equation.

It can be shown that the field anticommutation relation in terms of the
operators $a,b$ is
$$\left\{\psi_s(t,x),\psi_r^\dagger(t,y)\right\}=\delta^3(\mathbf{x-y})\delta_{rs}
\iff
\left\{a_s(\mathbf{p}),a_r^\dagger(\mathbf{p'})\right\}
=\left\{b_s(\mathbf{p}),b_r^\dagger(\mathbf{p'})\right\}=\delta^3(\mathbf{p-p'})
\delta_{rs}
$$
and the other are zero.

From this expansion of the field, $\psi^\dagger$ and $\dot\psi$ can be built
and the Hamiltonian turns out to be
$$H = \int \dd^3x \mathcal{H} = i\int \dd^3x\psi^\dagger\dot\psi$$
the strategy to integrate is to integrate over $\mathbf{x}$, yielding 
$\delta^3(\mathbf{p-p'})$ and then using this to apply the orthogonality relations
of the spinors $u,v,u^\dagger,v^\dagger$:
$$H = \int \dd^3p \sum_{s,r = 1}^{2} \frac{1}{2}\left[u_s^\dagger(\mathbf{p})
u_r(\mathbf{p})a_s^\dagger(\mathbf{p})a_r(\mathbf{p}) - 
v_s^\dagger(\mathbf{p})v_s(\mathbf{p})v_r(\mathbf{p})b_s(\mathbf{p})b_r^\dagger(
\mathbf{p})\right]$$
since
$$u^\dagger_s u_r =  v_s^\dagger v_r=2E_p\delta_{rs},$$and
$$ b_sb_s^\dagger = -b_s^\dagger b_s + \delta^3(0), $$
after taming the infinite offset, the Hamiltonian finally is
$$N(H) = \int \dd^3p E_p \sum_{s}^{}\left[a_s^\dagger(\mathbf{p})a_s(\mathbf{p})
+b_s^\dagger(\mathbf{p})b_s(\mathbf{p})\right].$$
\textbf{Note} Had the relations between $\psi,\pi_\psi$ been commutators,
the sign of the $b^\dagger_sb_s$ would have changed, giving way to a theory
without \textit{definite positive} energy, which is non-physical.

The operators have a similar physical interpretation: $a_s^\dagger(\mathbf{p})$
creates a fermion with spin $s$ and momentum $\mathbf{p}$, while
$b^\dagger_s(\mathbf{p})$ creates an antifermion with same spin and momentum.

\textbf{Note} The form of the Hamiltonian is remarkable: the terms
$a_s^\dagger a_s$, $b_s^\dagger b_s$ are number operators, so the total energy 
$H$ is the energy $E_p$ multiplied by the total number of particles and
antiparticles of momentum $p^i$ and integrated over all the possible values of
$p^i$.

\textbf{Note} This formalization has a \textbf{built-in Pauli exclusion principle,
} since applying two creation operators with same quantum number $s$ and same
momentum $p$ to the vacuum yields
$$a^\dagger_s(p)a^\dagger_s(p)\ket{0} = -a^\dagger_s(p)a^\dagger_s(p)\ket{0}
\equiv 0$$
while, if $r\neq s$ or $p\neq p'$,
$$a^\dagger_s(p)a^\dagger_r(p')\ket{0} = -a^\dagger_r(p')a^\dagger_s(p)\ket{0}
\neq0$$
since $\left\{a_s^\dagger(\mathbf{p}),a_r^\dagger(\mathbf{p'})\right\}=0$: this
closely resembles the Fermi-Dirac Statistic and is a direct consequence of the
choice of anticommutators over commutators.

Using again normal ordering and the F.Ts of the fields, the momentum operator
is
$$P^i = \int \dd^3x\psi^\dagger(-i\pd^i)\psi\to
N(P^i) = \int \dd^3pp^i\sum_{s}^{}(a^\dagger_sa+b^\dagger_sb_s).$$
Again this is the number of particles and antiparticles of momentum
$p^i$, multiplied by their momentum and integrated over all possible
values of it.
\subsubsection{Noether current and charges}
It was shown that the phase transformation
$$\psi\to \psi'=e^{i\alpha}\psi$$
is a symmetry for the Lagrangian and its associated Noether current is
$$J^\mu = \bar\psi\gamma^\mu\psi$$
with charge
$$Q = \int \dd^3x J_0 = \int \dd ^3x \psi^\dagger\psi$$
replacing the F.Ts into the integral (with the same \textit{caveat} of integrating
first over $x$) yields
$$Q = \int \dd^3p \sum_{s}^{}a^\dagger_sa_s + b_sb_s^\dagger$$
the normal ordering yields
$$N(Q) = \int \dd^3p \sum_{s}^{}\left[a^\dagger_sa_s - b_s^\dagger b_s\right]$$
this implies that the number of fermions - the number of antifermions \textbf{is
conserved}: so every time a fermion is created, an antifermion must also be
created [ragionamento su H e P, perso...].

Using Noether's theorem on Lorentz invariance of the Dirac field, the
Energy-momentum tensor turns out to be
$$T^{\mu\nu} = i\bar\psi\gamma^\mu\pd^\nu\psi$$
and the angular momentum becomes
$$M^{\lambda\mu\nu} = ... = i\bar\psi\gamma^\lambda(\underbrace{x^\mu\pd^\nu -
x^\nu\pd^\mu}_{\text{orbital ang mom}} \underbrace{-\frac{i}{2}\sigma^
{\mu\nu}}_\text{spin})$$
[RIVEDI!]
\section{The Propagator}
Until now, time was fixed, so the next natural step is to consider the
(space-)time evolution of the fields.

The propagator in QM is defined as
$$\psi(\mathbf{x},t) = \int \dd^3y K(\mathbf{y},t',\mathbf{x},t)\psi
(\mathbf{y},t')$$

\textbf{Def.} The free propagator of a field $\phi$ is the function
$$G(x^\mu,y^\nu) = \vartheta(x^0-y^0)\bra{0}\phi(x^\mu)\phi^\dagger(y^\nu)\ket{0}
+\vartheta(y^0-x^0) \bra{0}\phi(y^\mu)\phi^\dagger(x^\mu)\ket{0}$$
The first sandwich part can be interpreted as follows: $\phi^\dagger$ creates a
probe particle at spacetime $y^\nu$, the system is then left to evolve, until at
the spacetime point $x^\mu$ the particle is destroyed by $\phi$; the Heaviside
theta ensures the particle is destroyed only after it is created. The second
addend is the exact same, but with the creation and destruction of an
antiparticle.

The propagator is the \textbf{transition probability amplitude} of returning to
the ground state, which, in the case of free theories, is the vacuum $\ket{0}$.

